\section{De la llamada a un LLM a los sistemas Agent}

\subsection{El ciclo mínimo del Agent}

Los límites de entrada y salida de una sola llamada a un LLM son fijos; un Agent, en cambio, mantiene un objetivo y un estado dentro de un ciclo: observa el entorno, decide una acción, ejecuta una herramienta, lee la retroalimentación y luego decide el siguiente paso. Su abstracción mínima es:

$$
s_{t+1}=T(s_t,a_t),\qquad a_t\sim\pi_{\theta}(\cdot\mid s_t,g).
$$

\begin{itemize}
    \item $g$: objetivo del usuario;
    \item $s_t$: estado visible en el paso $t$, que incluye el historial, los resultados de las herramientas y el entorno externo;
    \item $a_t$: la acción elegida por el modelo;
    \item $T$: la transición de estado tras la ejecución de la acción por el entorno;
    \item $\pi_{\theta}$: la política del modelo que decide la siguiente acción.
\end{itemize}

\begin{figure}[H]
\centering
\includegraphics[width=0.9\textwidth]{figures/agent-loop.jpg}
\caption{El Agent forma un ciclo cerrado mediante acciones y retroalimentación del entorno.\protect\footnotemark}
\end{figure}
\footnotetext{Intervalo de la explicación en video: 00:42:20--00:43:53.}

\begin{importantbox}{El núcleo de un Agent no es "llamar al modelo unas cuantas veces más"}
Solo cuando las decisiones posteriores leen resultados de ejecución reales y pueden modificar el plan a partir de ellos, el sistema forma un ciclo de retroalimentación cerrado. La generación de texto en varias rondas sin estado del entorno, ejecución de herramientas ni consumo de retroalimentación se parece más a un prompt pipeline que a un Agent completo.
\end{importantbox}

\subsection{Agentic Workflow: escribir estructura confiable en código}

El agentic workflow usa un flujo de control predefinido para organizar los modelos y las herramientas. Sacrifica parte de la autonomía a cambio de previsibilidad, observabilidad y facilidad de depuración. Por ejemplo, en un flujo de generar, luego revisar y reintentar si falla, la lógica de control está definida por el programa, no decidida libremente por el modelo cada vez.

\begin{figure}[H]
\centering
\includegraphics[width=0.94\textwidth]{figures/agentic-workflows.jpg}
\caption{Dos tipos de workflow: generación--evaluación en serie y generación--agregación en paralelo.\protect\footnotemark}
\end{figure}
\footnotetext{Intervalo de la explicación en video: 00:43:53--00:44:31.}

El curso enumera cinco patrones comunes:

\begin{table}[H]
\centering
\small
\begin{tabular}{p{3.2cm}p{5.4cm}p{5.2cm}}
\toprule
\textbf{Patrón} & \textbf{Mecanismo} & \textbf{Escenario de aplicación} \\
\midrule
Prompt chaining & La salida del paso anterior se vuelve la entrada del siguiente & Generación de documentos, extracción de información, transformación por etapas \\
Routing & Primero se determina el tipo de solicitud y luego se asigna a un modelo o herramienta especializados & Distribución de atención al cliente, cascada de modelos, clasificación de riesgo \\
Parallelization & Se generan varios candidatos o se procesan varias subtareas al mismo tiempo & Muestreo diverso, recuperación independiente, análisis por lotes \\
Orchestrator--worker & Un planificador central divide la tarea, varios workers la ejecutan y al final se integra & Modificación de bases de código, investigación profunda, informes complejos \\
Evaluator--optimizer & El generador y el revisor iteran hasta cumplir el criterio o agotar el presupuesto & Revisión de escritura, pruebas de código, optimización de soluciones \\
\bottomrule
\end{tabular}
\caption{Patrones comunes de agentic workflow}
\end{table}

\begin{figure}[H]
\centering
\includegraphics[width=0.94\textwidth]{figures/workflow-patterns.jpg}
\caption{Un workflow se compone de llamadas a LLM, herramientas, búsqueda, verifier, critic y judge.\protect\footnotemark}
\end{figure}
\footnotetext{Intervalo de la explicación en video: 00:44:32--00:49:00.}

\subsection{Compromisos de ingeniería entre workflow y Agent autónomo}

\begin{itemize}
    \item \textbf{Workflow fijo}: alta controlabilidad, fácil de probar, límite de costo claro, pero carece de adaptabilidad ante situaciones nuevas;
    \item \textbf{Agent autónomo}: puede planificar dinámicamente y elegir herramientas, pero la explosión de ramas de comportamiento hace que los errores se acumulen en trayectorias largas;
    \item \textbf{Sistema híbrido}: fija de forma rígida los límites de alto riesgo, los permisos y los criterios de aceptación, y deja la búsqueda abierta al modelo.
\end{itemize}

\begin{warningbox}{A mayor autonomía, menos basta con evaluar solo la respuesta final}
Un sistema de trayectoria larga puede llegar por accidente a un resultado correcto y aun así haber pasado por acciones peligrosas, ciclos ineficaces o costos inaceptables. La evaluación debe registrar a la vez la calidad final, la corrección de las llamadas a herramientas, el número de pasos, la latencia, el costo, los límites de permisos y el comportamiento de recuperación ante fallos.
\end{warningbox}

\subsection{Resumen de la sección}

El Agent extiende el razonamiento hasta convertirlo en un proceso de acción con estado; el agentic workflow escribe en el flujo de control los patrones comunes de descomposición, routing, paralelización y revisión. Ambos requieren los temas posteriores del curso: la planificación decide qué hacer, las herramientas proporcionan retroalimentación del entorno, el verifier determina si es correcto, y el entrenamiento y la memoria hacen que el sistema mejore la próxima vez.
